\section*{第 \ref{ch:g11:electric-gravitational-fields} 章——电场与引力场}

\begin{solution}{exo:g11:electric-gravitational-fields:1}
$E = \num{5.0e-3} / \num{2.0e-6} = \qty{2.5e3}{N/C}$，方向朝北。
在 $q'$ 上：$F = \num{4.0e-6} \times \num{2.5e3} = \qty{1.0e-2}{N}$，
方向朝\emph{南}（$q' < 0$）。
\end{solution}

\begin{solution}{exo:g11:electric-gravitational-fields:2}
\emph{1.} $E = \num{8.99e9} \times \num{5.0e-9} / 0.30^2
\approx \qty{5.0e2}{N/C}$，方向背离 $O$（$Q > 0$）。

\emph{2.} 要使 $k\abs{Q}/d^2$ 减半，需将 $d^2$ 加倍：
$d' = 0.30\sqrt{2} \approx \qty{0.42}{m}$。
\end{solution}

\begin{solution}{exo:g11:electric-gravitational-fields:3}
$E = 600 / \num{3.0e-3} = \qty{2.0e5}{V/m}$；
$F = eE = \num{1.60e-19} \times \num{2.0e5} = \qty{3.2e-14}{N}$。
\end{solution}

\begin{solution}{exo:g11:electric-gravitational-fields:4}
$g_{\text{火星}} = P/m = 930/250 = \qty{3.7}{N/kg}$。宇航员：
$P = 80 \times 3.72 \approx \qty{3.0e2}{N}$（地球上
约为 $\qty{780}{N}$）。
\end{solution}

\begin{solution}{exo:g11:electric-gravitational-fields:5}
（a）错误：每一点的场只有一个方向，因此电场线永不
相交。（b）错误：电场线从 $+$ 出发，到 $-$ 终止。（c）正确：电场线的疏密
表示场的强弱。
\end{solution}

\begin{solution}{exo:g11:electric-gravitational-fields:6}
\emph{1.} $d = \num{6.37e6} + \num{4.2e5} = \qty{6.79e6}{m}$，因此
$g = \num{6.67e-11} \times \num{5.97e24} / (\num{6.79e6})^2
\approx \qty{8.6}{N/kg}$——为地表值的 $88\%$。

\emph{2.} 二者漂浮，是因为空间站与宇航员一起
处于自由落体状态，而并非因为重力不存在。
\end{solution}

\begin{solution}{exo:g11:electric-gravitational-fields:7}
$E = m_e g / e = \num{9.11e-31} \times 9.81 / \num{1.60e-19}
\approx \qty{5.6e-11}{V/m}$：约比示波器中的
电场弱 $\num{4e14}$ 倍。在电子物理中，重力完全可以忽略。
\end{solution}

\begin{solution}{exo:g11:electric-gravitational-fields:8}
\emph{1.} $E = \num{5.0e3} / \num{2.0e-2} = \qty{2.5e5}{V/m}$；静止时
$qE = mg$，因此
$q = \num{2.0e-6} \times 9.81 / \num{2.5e5} \approx \qty{7.8e-11}{C}$。

\emph{2.} 正颗粒所受电场力必须竖直向上，因此 $\vect E$
方向竖直向上：\emph{下}极板带正电。

\emph{3.} $n = \num{7.85e-11} / \num{1.60e-19} \approx \num{4.9e8}$
个元电荷。
\end{solution}

\begin{solution}{exo:g11:electric-gravitational-fields:9}
在线段外侧，两个场方向相同；在线段之间，二者
方向相反——要使合场为零，该点应更靠近较弱的源 $+Q$。
设 $x$ 为到 $+Q$ 的距离，则 $kQ/x^2 = 4kQ/(0.30 - x)^2$，
故 $0.30 - x = 2x$，于是 $x = \qty{0.10}{m}$，即距 $+Q$ 处。
\end{solution}

\begin{solution}{exo:g11:electric-gravitational-fields:10}
$F_e = ke^2/d^2 = \num{8.99e9} \times (\num{1.60e-19})^2 /
(\num{5.3e-11})^2 \approx \qty{8.2e-8}{N}$；
$F_g = G m_p m_e / d^2 = \num{6.67e-11} \times \num{1.67e-27} \times
\num{9.11e-31} / (\num{5.3e-11})^2 \approx \qty{3.6e-47}{N}$。比值
$F_e / F_g \approx \num{2.3e39}$。
\end{solution}

\begin{solution}{exo:g11:electric-gravitational-fields:11}
$E = \sqrt{300^2 + 400^2} = \qty{500}{N/C}$，方向北偏东
$\tan^{-1}(400/300) \approx \ang{53}$。对 $q < 0$：
$F = \num{2.0e-6} \times 500 = \qty{1.0e-3}{N}$，方向与场相反——
$\ang{53}$ 南偏西。
\end{solution}

\begin{solution}{exo:g11:electric-gravitational-fields:12}
$GM/(R+h)^2 = \tfrac12\,GM/R^2$ 得 $R + h = R\sqrt{2}$，因此
$h = (\sqrt{2} - 1)R \approx \qty{2.6e6}{m}$（约 \qty{2600}{km}）。
当 $h = R$ 时，$d = 2R$，重力加速度为地表值的四分之一。
\end{solution}

\begin{solution}{exo:g11:electric-gravitational-fields:13}
\emph{1.} $AM = \sqrt{3.0^2 + 4.0^2} = \qty{5.0}{cm}$；每个电荷
产生的场为 $E_1 = \num{8.99e9} \times \num{2.0e-9} / 0.050^2
\approx \qty{7.2e3}{N/C}$。

\emph{2.} 由对称性，沿 $AB$ 方向的分量相互抵消；每个场
沿角平分线、背离线段方向的分量为 $E_1 \times 4/5$，
故 $E = 2 \times \num{7.19e3} \times 0.80 \approx \qty{1.2e4}{N/C}$，
方向沿角平分线，背离 $[AB]$。
\end{solution}

\begin{solution}{exo:g11:electric-gravitational-fields:14}
盒内的自由电荷在外电场作用下移动：$+$ 电荷聚集在
一个面上，$-$ 电荷聚集在另一面上，这些移动的电荷产生
与外电场方向相反的内电场。它们持续移动，直至
盒内总电场恰好为零，此时达到平衡。引力
只有一种符号的源：每个质量都\emph{增加}一个吸引场，
没有任何质量能够抵消它。引力屏蔽需要负质量，
而负质量并不存在。
\end{solution}

\begin{solution}{exo:g11:electric-gravitational-fields:15}
\emph{1.} 重力 $m\vect g$（竖直向下），张力 $\vect T$（沿
细线），电场力 $q\vect E$（水平）。静止时，
$T\sin\theta = qE$，$T\cos\theta = mg$；两式相除：
$\tan\theta = qE/(mg)$。

\emph{2.} $\tan\theta = \num{2.0e-7} \times \num{1.0e4} /
(\num{5.0e-4} \times 9.81) \approx 0.41$，故 $\theta \approx \ang{22}$；
$T = mg/\cos\theta \approx \qty{5.3e-3}{N}$.
\end{solution}

\begin{solution}{pb:g11:electric-gravitational-fields:1}
\textbf{1.} 匀强电场，垂直于极板，方向向下（从
上方 $+$ 极板指向下方 $-$ 极板）；场强处处相同，
即液滴无论漂移到何处，所受的力都相同。
\textbf{2.} $E = U/d = 800 / \num{8.0e-3} = \qty{1.0e5}{V/m}$.
\textbf{3.} 电荷为负，故电场力方向
\emph{与 $\vect E$ 相反}：竖直向上。仅当 $+$ 极板在上方时，$\vect E$ 才
竖直向下，电场力与重力相抗衡。
\textbf{4.} $P = mg = \num{4.90e-15} \times 9.81 \approx
\qty{4.8e-14}{N}$.
\textbf{5.} $\abs{q}E = mg$.
\textbf{6.} $\abs{q} = mg/E = mgd/U$.
\textbf{7.} $\abs{q} = \num{4.81e-14} \times \num{8.0e-3} / 800
\approx \qty{4.8e-19}{C}$.
\textbf{8.} $n = \num{4.8e-19} / \num{1.60e-19} = 3$ 个多余电子。
\textbf{9.} $E$ 增大，故 $\abs{q}E > mg$：合力向上。
\textbf{10.} 由 $\abs{q}\,U/d = mg$，得 $U = mgd/\abs{q}$——
与电荷成反比。
\textbf{11.} $\abs{q} = mgd/U = \num{3.85e-16}/1200 \approx
\qty{3.2e-19}{C} = 2e$：比 $3e$ 减少，液滴\emph{失去}了一个
电子。
\textbf{12.} $\abs{q} = \num{3.85e-16}/U$: \qty{3.2e-19}{C},
\qty{4.8e-19}{C}, \qty{6.4e-19}{C}, \qty{8.0e-19}{C},
\qty{9.6e-19}{C}.
\textbf{13.} 它们是 $2e, 3e, 4e, 5e, 6e$，其中
$e = \qty{1.60e-19}{C}$.
\textbf{14.} 每一步使 $\abs{q}$ 恰好改变
\qty{1.6e-19}{C}：每次只有一个电子附着上去。
\textbf{15.} 此时所需电荷为 $\abs{q} = \num{3.85e-16}/700 \approx
\qty{5.5e-19}{C} = 3.4\,e$——不是整数个电子，因此
没有任何辐射爆发能够产生它。
\textbf{16.} 电荷是\emph{量子化}的，只能取
元电荷 $e = \qty{1.60e-19}{C}$ 的整数倍——这正是
密立根所测得的量。
\textbf{17.} $q \leftrightarrow m$，$E \leftrightarrow g$，以及
$\abs{q}E \leftrightarrow mg$：正是重力本身。
\textbf{18.} $g_{\text{球体}} = GM/d^2 = \num{6.67e-11} \times 1000 /
1.0^2 = \qty{6.7e-8}{N/kg}$；液滴所受的力
$\num{4.90e-15} \times \num{6.7e-8} \approx \qty{3.3e-22}{N}$——比其重力小约
$\num{1.5e8}$ 倍。毫无希望。
\textbf{19.} 由于 $G$ 很小，只有头顶上方行星大小的质量
才能与地球的引力相抗衡——而没有任何质量能够排斥。巨大的 $k$
以及两种电荷符号，使台面上的一块极板
在任一方向上都能超过行星的引力。
\textbf{20.} 一个电子：$\Delta\abs{q}/\abs{q} = e/3e \approx 33\%$；
一个油分子：$\Delta m/m = \num{5e-25} / \num{4.90e-15}
\approx \num{e-10}$。电荷阶梯使平衡电压
改变数百伏特；质量阶梯在任何电压下都无法察觉。
因此，该天平能够分辨“电的原子”——
$e = \qty{1.60e-19}{C}$——而油分子则模糊成一个
连续体。
\end{solution}
